\documentclass[11pt]{article}

\usepackage[margin=1in]{geometry}
\usepackage{cmap}
\usepackage[T1]{fontenc}
\usepackage[utf8]{inputenc}
\usepackage{booktabs}
\usepackage{longtable}
\usepackage{array}
\usepackage{xcolor}
\usepackage{hyperref}
\usepackage{listings}
\usepackage{amsmath}

\hypersetup{
  hidelinks,
  pdftitle={MATRIX First Release Manual},
  pdfauthor={MATRIX Development Team}
}

\definecolor{MatrixBlue}{HTML}{245C73}
\definecolor{MatrixSlate}{HTML}{4A5561}
\definecolor{MatrixGold}{HTML}{C88A2D}
\definecolor{MatrixBg}{HTML}{F7F8F6}
\emergencystretch=2em

\lstdefinestyle{matrix}{
  basicstyle=\ttfamily\small,
  breaklines=true,
  columns=fullflexible,
  keepspaces=true,
  frame=single,
  rulecolor=\color{MatrixSlate!45},
  backgroundcolor=\color{MatrixBg},
  xleftmargin=0.5em,
  xrightmargin=0.5em
}

\newcommand{\framework}{\texttt{MATRIX}}
\newcommand{\oracle}{\texttt{ORACLE}}
\newcommand{\xyzin}{\texttt{xyzin}}
\newcommand{\sectionname}[1]{\texttt{\##1}}

\title{\textbf{MATRIX First Release Manual}\\
\large Operational guide for the implemented ORACLE/MATRIX tools}
\author{MATRIX Development Notes}
\date{\today}

\begin{document}
\maketitle

\begin{abstract}
This manual defines the first release boundary of \framework{} and gives the
operational contract for the implemented tools.  \framework{} is the framework:
each tool reads and writes shared enriched \xyzin{} sections, uses shared
libraries for repeated tasks, and can run independently when its required
sections are already present.  Fragment search, nano-LEGO assembly and the
production TRINITY optimizer are reserved for the next release; this release is
closed by the preprocessing, GIC, GF/PED, MORPHEUS, rovibrational, anharmonic,
DVR, QM-adapter and GUI workflows documented here.
\end{abstract}

\tableofcontents
\newpage

\section{Release Boundary}

The first \framework{} release is a refactored, documented and testable
replacement for the corresponding implemented parts of Merlino/ORACLE.  The
release is complete when the following conditions hold:

\begin{itemize}
  \item every implemented scientific tool can run from the command line using
        only its declared \xyzin{} sections;
  \item the GUI launches the same public commands and does not contain private
        parsers or scientific algorithms;
  \item all external file formats enter through one adapter per format;
  \item the same libraries are used for geometry, topology, symmetry,
        normal-mode data, diagonalization and publication exports;
  \item manuals identify consumed sections, produced sections, standalone
        commands, GUI placement and regression tests.
\end{itemize}

The next release owns robust fragment search, nano-LEGO fragmentation and
assembly, LCB25 fragment indexing, the production TRINITY external
energy/gradient optimization loop, PySCF launch/analysis adapters, VSCF for
local anharmonic oscillators, vibronic spectroscopy, ionization potentials,
electron affinities and online help/manual generation.  The current release
keeps the fragment and TRINITY schemas, command skeletons and GUI entry points
so downstream design does not need another reorganization.

\section{Shared Container Contract}

The central object is the enriched \xyzin{} file.  Each tool appends or updates
sections instead of passing private scratch formats between modules.  A tool
must not reparse Gaussian, Molpro, MRCC, SMILES or geometry data if the
corresponding shared section is already present.

\begin{longtable}{
  >{\raggedright\arraybackslash}p{0.25\linewidth}
  >{\raggedright\arraybackslash}p{0.33\linewidth}
  >{\raggedright\arraybackslash}p{0.32\linewidth}}
\toprule
Section & Owner & Main consumers \\
\midrule
\sectionname{BASIC} & ORACLE-Babel / chemistry layer & all molecular tools \\
\sectionname{SYMMETRY} & ORACLE-Babel & NEO/GICForge, GUI \\
\sectionname{TOPOLOGY} & ORACLE-Babel & synthons, fragments, GICs \\
\sectionname{SYNTHONS} & ORACLE-Babel & GF/PED, MORPHEUS, future fragments \\
\sectionname{GIC}, \sectionname{SYCART} & NEO/GICForge & GF/PED, MORPHEUS, TRINITY \\
\sectionname{CARTESIAN\_HESSIAN} & QM adapters & GF/PED, diagnostics \\
\sectionname{NORMAL\_MODES} & QM adapters & vibrational comparison, hybrid spectra \\
\sectionname{QFF} & QM adapters & VPT2/VCI \\
\sectionname{ROTATIONAL} & rovib/QM adapters & WMS-Rot, thermo \\
\sectionname{VIBRATIONAL} & rovib/QM adapters & spectra, thermo, DOS \\
\sectionname{GF\_PED} & GF/PED & GUI, vibrational workbench \\
\sectionname{MORPHEUS} & MORPHEUS/SEFit & GUI, reports \\
\sectionname{THERMO} & thermo & publication exports \\
\sectionname{VPT2\_VCI} & VPT2/VCI & vibrational workbench \\
\sectionname{DVR} & DVR & GUI, post-run analysis \\
\bottomrule
\end{longtable}

The practical rule is simple: if two tools need the same information, that
information belongs in a shared section and in a shared library, not in two
tool-local parsers.

\section{ORACLE-Babel and Preprocessing}

ORACLE-Babel is the preprocessing gate.  It obtains an initial Cartesian
geometry by reading an XYZ file directly, importing from a supported QM format,
or building from SMILES through RDKit.  It then performs symmetry detection,
topology analysis, synthon typing and validation once, and writes the results
back to the same \xyzin{} container.

\begin{lstlisting}[style=matrix]
oracle babel preprocess input.xyz molecule.xyzin
oracle babel preprocess smiles.txt molecule.xyzin --source-format smiles
oracle validate molecule.xyzin
\end{lstlisting}

The output sections are normally \sectionname{SOURCE}, \sectionname{BASIC},
\sectionname{SYMMETRY}, \sectionname{TOPOLOGY} and \sectionname{SYNTHONS}.
Downstream tools must consume those sections rather than re-detecting bonds,
rings, point groups or atom classes.

Avogadro remains the preferred manual 3D inspection and editing tool.  After
the user accepts the molecular structure, the GUI displays the 2D molecular
view and all later tasks work from the enriched \xyzin{} state.

\section{QM Adapters}

QM adapters are the only code allowed to understand external program outputs.
Gaussian, Molpro and MRCC data are promoted into shared sections.  Scientific
tools consume these sections and never scan log files privately.

\begin{lstlisting}[style=matrix]
oracle gaussian promote-fchk calc.fchk molecule.xyzin
oracle gaussian promote-rovib calc.log molecule.xyzin
oracle gaussian promote-electronic calc.log molecule.xyzin
oracle molpro promote-output molpro.out molecule.xyzin
oracle mrcc promote-output mrcc.out molecule.xyzin
\end{lstlisting}

The main promoted sections are \sectionname{CARTESIAN\_HESSIAN},
\sectionname{NORMAL\_MODES}, \sectionname{QFF}, \sectionname{ROTATIONAL},
\sectionname{VIBRATIONAL}, \sectionname{DELTABVIB},
\sectionname{ELECTRONIC}, \sectionname{TRANSITIONS} and
\sectionname{ORBITALS}.  Orbital files are passed to Molden, Avogadro or
MOrbVis-compatible viewers through the GUI command layer.

\section{NEO / GICForge}

The implemented GIC tool is still exposed as \texttt{gicforge} in the command
line, with \texttt{neo} as an alias.  The planned tool name is NEO,
``Nonredundant Equivariant Orthogonalizer''.  The specialist manual
\texttt{gicforge\_manual.pdf} remains the detailed reference for primitive
construction, ring puckering, symmetry projectors, special fragment-aware
coordinates, B-matrix construction and Python/Fortran parity.

\begin{lstlisting}[style=matrix]
oracle gicforge molecule.xyzin
oracle neo molecule.xyzin
oracle gic-bmatrix oracle.gic.definition.json geometry.xyz
\end{lstlisting}

NEO consumes \sectionname{VALIDATION}, \sectionname{SYMMETRY},
\sectionname{TOPOLOGY}, \sectionname{SYNTHONS} and optional
\sectionname{FRAGMENTS}.  It produces \sectionname{GIC} and
\sectionname{SYCART}.  Symmetry-adapted GIC labels begin with the irrep label;
totally symmetric blocks can therefore be selected directly in optimization or
least-squares workflows.

\section{GF/PED}

GF/PED consumes a frozen GIC basis and a Cartesian Hessian.  The specialist
manual \texttt{gf\_manual.pdf} is the detailed reference for Wilson GF theory,
internal force constants, PEDs, local force-field approximations and
Pulay-style scaling of force constants in GIC space.

\begin{lstlisting}[style=matrix]
oracle gf --xyzin molecule.xyzin --out gf.report
oracle gf --xyzin molecule.xyzin --csv-dir gf_tables
\end{lstlisting}

The module writes \sectionname{GF\_PED}.  It can subtract electrostatic and van
der Waals non-bonded terms before a local force-field model is built.  Frequency
scaling is not a MATRIX vibrational-spectrum operation; scaling belongs here,
at the force-constant level, where diagonal class scaling and geometric-mean
off-diagonal scaling are chemically meaningful.

\section{MORPHEUS / SEFit}

MORPHEUS performs semiexperimental geometry refinement.  It includes the SEFit
single-structure workflow and the multi-structure class-correction workflows.
The detailed manuals are \texttt{morpheus\_manual.pdf} and
\texttt{multistructure\_manual.pdf}.

\begin{lstlisting}[style=matrix]
oracle semiexp --xyzin molecule.xyzin --job job.toml --outdir run
oracle semiexp-ensemble ensemble.toml --outdir ensemble_run
oracle semiexp-benchmark snapshot.json --outdir benchmark
\end{lstlisting}

MORPHEUS consumes isotopologue data and optional \sectionname{GIC},
\sectionname{SYCART} or \sectionname{MORPHEUS} state.  It writes
\sectionname{MORPHEUS}, diagnostics, CSV tables and paper-oriented summaries.
The local \texttt{se\_geometries} library is an implemented support library;
large-scale fragment search and nano-LEGO assembly remain next-release work.

\section{Rotational Spectroscopy / WMS-Rot}

The rotational workbench uses the local WMS-Rot Hamiltonian engine imported
from the group web service.  MATRIX does not reimplement the Hamiltonian.  It
normalizes constants, dipoles and vibration-rotation corrections through
\sectionname{ROTATIONAL}, \sectionname{VIBRATIONAL} and
\sectionname{DELTABVIB}, then calls the local engine.

\begin{lstlisting}[style=matrix]
oracle rovib summarize molecule.xyzin
oracle rovib wmsrot-input molecule.xyzin --out wmsrot.in
oracle rovib wmsrot-run molecule.xyzin --out rotational_lines.csv
\end{lstlisting}

The run writes line lists suitable for plotting and publication export.  Future
database comparison should query microwave/submillimeter catalogs through a
shared experimental-spectrum adapter, not from GUI-local code.  The planned
catalog targets are CDMS, JPL and Splatalogue, with molecule identifiers,
frequency windows, isotopic species and provenance saved in export metadata.

\section{Vibrational Spectroscopy}

The vibrational workbench plots IR, Raman, VCD and ROA spectra when the
corresponding harmonic or anharmonic data are present in
\sectionname{VIBRATIONAL}.  It can compare spectra from one or two \xyzin{}
files.  For ordinary IR and Raman comparisons, the second trace is plotted with
negative intensity.  VCD and ROA are signed observables and are therefore not
mirrored.

\begin{lstlisting}[style=matrix]
oracle rovib vib-spectrum molecule.xyzin --observable IR \
  --source harmonic --csv ir.csv --plot ir.svg

oracle rovib vib-compare level1.xyzin level2.xyzin \
  --observable IR --first-source harmonic --second-source anharmonic \
  --csv compare.csv --plot compare.svg
\end{lstlisting}

Hybrid spectra are implemented as
\[
  \nu_\mathrm{hybrid}
  =
  \nu_\mathrm{harm}(\mathrm{level1})
  +
  \left[
    \nu_\mathrm{anh}(\mathrm{level2})
    -
    \nu_\mathrm{harm}(\mathrm{level2})
  \right].
\]
Before applying the correction, MATRIX reads \sectionname{NORMAL\_MODES} from
both files and performs a global assignment that maximizes the absolute
normal-mode overlap.  If any assigned overlap is below the threshold, the
hybrid spectrum is rejected.

\begin{lstlisting}[style=matrix]
oracle rovib vib-spectrum level1.xyzin --source hybrid \
  --level2-xyzin level2.xyzin --observable IR \
  --csv hybrid.csv --plot hybrid.svg \
  --mode-match-csv mode_matches.csv --min-mode-overlap 0.70
\end{lstlisting}

Experimental NIST IR spectra are fetched automatically only when the NIST
JCAMP record declares a gas-phase state.  Condensed-phase records, missing
records and unparseable records are returned to the GUI as states requiring
user instructions.

\begin{lstlisting}[style=matrix]
oracle rovib nist-ir 74-82-8 --out methane_nist_ir.csv
\end{lstlisting}

\section{Thermochemistry and Kinetics}

The implemented thermochemistry tool consumes normalized molecular, rotational
and vibrational sections and writes \sectionname{THERMO}.  The GUI
Thermo/Kinetics workbench displays the generated thermochemical tables and
rovibrational density-of-states outputs.  Production kinetics models and
\sectionname{KINETICS} are a future extension boundary, not a blocker for this
release.

\begin{lstlisting}[style=matrix]
oracle thermo molecule.xyzin --out thermo.report
oracle rovib dos molecule.xyzin --out dos_vib.dat
oracle rovib dos-rovib molecule.xyzin --out dos_rovib.dat --q-out rovib_qt.dat
\end{lstlisting}

Publication exports should keep temperature, pressure, symmetry number,
frequency set, rotational constants and data provenance in the table metadata.

\section{VPT2 / VCI}

The VPT2/VCI package consumes \sectionname{QFF}.  Gaussian FCHK or QFF text is
only an adapter input; standalone operation reads the normalized \xyzin{}
section.  Reports and CSV tables are attached to the vibrational workbench.

\begin{lstlisting}[style=matrix]
oracle vpt2-vci --xyzin molecule.xyzin --run-dir vpt2_run
oracle vpt2-vci --xyzin molecule.xyzin --run-dir vci_run --model vci
\end{lstlisting}

VPT2/VCI uses the shared diagonalization policy where matrix diagonalization is
the dominant cost.  GPU acceleration is delegated to the shared
\texttt{oracle\_core.diagonalizer} layer when supported by the local machine.

\section{DVR}

DVR workflows prepare and run scan/path DVR calculations and then normalize the
post-run outputs into the GUI state.  Large Hamiltonian diagonalizations are
routed through the same shared diagonalizer used by other MATRIX modules.

\begin{lstlisting}[style=matrix]
oracle dvr prepare molecule.xyzin --scan-log scan.log --run-dir dvr_run
oracle dvr run --xyzin molecule.xyzin --run-dir dvr_run
oracle dvr status molecule.xyzin
\end{lstlisting}

The implemented GUI state is \sectionname{DVR}.  New DVR engines must not add
private diagonalization wrappers; they should call the shared diagonalizer so
CPU/GPU policy is controlled in one place.

\section{ORACLE GUI}

The GUI keeps the ORACLE name in the MATRIX framework.  Its planned expansion
is ``Operator for Routing, Analysis, Control, Launch and Exploration''.  The
GUI owns project state, command construction, missing-section guidance, file
selection, external viewer launches and publication export orchestration.  It
does not own scientific parsers or algorithms.

\begin{lstlisting}[style=matrix]
python -m oracle_gui molecule.xyzin
oracle-gui molecule.xyzin
\end{lstlisting}

The first-release GUI exposes:

\begin{itemize}
  \item project dashboard and section readiness;
  \item structure/synthon/topology inspection and Avogadro handoff;
  \item NEO/GICForge, GF/PED, MORPHEUS, VPT2/VCI, DVR and TRINITY panels;
  \item rotational, vibrational, electronic and thermochemistry workbenches;
  \item orbital viewer launch through Molden, Avogadro or MOrbVis-compatible
        external viewers;
  \item missing-section suggestions that point to the owning adapter or tool.
\end{itemize}

\section{Regression and Release Tests}

The release gate is the full test suite plus focused smoke tests for workflows
that cross tool boundaries.  The minimum command is:

\begin{lstlisting}[style=matrix]
source scripts/oracle_env.sh
oracle-set
python -m ruff check packages tests
python -m pytest -q
\end{lstlisting}

Important cross-module checks include:

\begin{itemize}
  \item Python/Fortran NEO/GICForge primitive, GIC and B-row agreement;
  \item GF/PED writing and reading \sectionname{GF\_PED};
  \item MORPHEUS standalone and multi-structure runs;
  \item WMS-Rot local line-list generation;
  \item vibrational mirror comparison and hybrid level1+level2 normal-mode
        matching;
  \item VPT2/VCI and DVR diagonalizer routing;
  \item GUI command builders matching the public CLI.
\end{itemize}

\section{First Release Closure}

With this manual and the specialist manuals present, the first \framework{}
release can be declared closed when the repository is clean, the full tests
pass, and the manuals are generated from source.  The next release then starts
from a clear scope: robust fragments, nano-LEGO, LCB25 fragment indexing,
production TRINITY optimization, PySCF adapters, VSCF, vibronic spectroscopy,
IP/EA workflows and online help/manual generation.

\end{document}
